\section{Arbitrary sections and their local obstruction}
\label{sec:arbitrary}

The distinction between arbitrary sections and Selmer sections already
appears when ramification is allowed only at $3$. The construction below
uses Shafarevich's arithmetic freeness theorem and the free-source method
of Hoshi. We give the equivariant construction over $\Q$ explicitly.
The exchanged arithmetic basis also appears in
\citet[Theorem 12, p.~12]{Medvedovsky2016}; the elementary argument is
included to make the descent under complex conjugation transparent.

Put $P=P_{\{3\}}$ and $\Gamma=\Gamma_{\{3\}}$ in the notation of
Section~\ref{sec:local-arithmetic-source}. Complex conjugation supplies an involution $c$, so
\[
\Gamma=P\rtimes\langle c\rangle.
\]
The field $K=\Q(\zeta_3)$ contains $\mu_3$, has class number one, and
has just one prime above $3$. Thus
\citet[\S4, example 1 after Theorem 5, p.~82]{Shafarevich1963}
gives relation rank zero and generator rank
$[K:\Q]-\operatorname{rank}\mathcal O_K^\times=2$.
Consequently $P$ is free pro-$3$ of rank two; see also
\citet[Lemma 3.2]{Hoshi2010}.

\begin{lemma}[An exchanged free basis]\label{lem:exchanged}
There is a free pro-$3$ basis $g,h$ of $P$ with $h=cgc$.
Thus $c$ exchanges the two generators exactly.
\end{lemma}
\begin{proof}
The extensions $K(\zeta_9)$ and $K(\sqrt[3]{3})$ are independent
cyclic cubic extensions unramified outside $3$. Their Kummer classes
are $\zeta_3$ and $3$: the first has zero valuation, whereas the
second has $\lambda$-valuation two, nonzero modulo three. They
therefore give the two independent quotients of $P/\Phi(P)$.
Complex conjugation acts trivially on the cyclotomic quotient,
since $\Gal(\Q(\zeta_9)/\Q)$ is abelian, and by inversion on the
other quotient, whose extension over $\Q$ has group $S_3$.

Choose $g\in P$ whose Frattini vector has both eigencoordinates
nonzero, say $(a,b)$. The vector of $h=cgc$ is $(a,-b)$.
Their determinant is $-2ab\ne0$ in $\F_3$, so Burnside's basis
theorem says $g,h$ generate $P$. A generating pair of a rank-two
free pro-$3$ group is a free basis: the resulting surjective
endomorphism is invertible by the Hopfian property of finitely
generated profinite groups. Since $c^2=1$, it also sends $h$ to $g$.
\end{proof}

\begin{lemma}[Direct full-group cocycles]\label{lem:directcocycles}
For every pro-$3$ group $M$ with a continuous $\Gamma$-action and
every $a\in M$, there is a continuous $M$-cocycle $z_a$ with
\[
z_a(c)=1,\qquad z_a(g)=a,\qquad z_a(h)=c(a).
\]
\end{lemma}
\begin{proof}
The universal property of the free pro-$3$ basis gives a continuous
homomorphism
\[
P\longrightarrow M\rtimes P,\qquad
g\longmapsto(a,g),\quad h\longmapsto(c(a),h).
\]
The target is pro-$3$ and its projection to $P$ is the identity
on the free generators, hence everywhere. Conjugation by $(1,c)$
exchanges the two specified images. The homomorphism is therefore
$c$-equivariant and extends to a section
$\Gamma\to M\rtimes\Gamma$ sending $c$ to $(1,c)$. Its relative
cocycle is the asserted $z_a$. This constructs an actual full-group
cocycle, with no passage from graded lifts to an inverse limit.
\end{proof}

Let $E$ be the proper cubic and put $T=T_3E$, with the action supplied
by a standard cusp tangent.
\begin{lemma}[The compact cohomology]\label{lem:compactrank}
There is an isomorphism of compact $\Z_3$-modules
\[
H^1(\Gamma,T)=T/(g-1)T^{c=1}.
\]
This group is finitely generated and has rank one.
\end{lemma}
\begin{proof}
Since multiplication by two is invertible on $T$, every cocycle is
cohomologous to one whose value at $c$ is zero. Such a normalized
cocycle is determined by its arbitrary value $x$ at $g$, with value
$c(x)$ at $h$. Conversely every $x\in T$ occurs by
Lemma~\ref{lem:directcocycles}. A coboundary preserves the
normalization exactly when it comes from $y\in T^{c=1}$, and its
value at $g$ is $(g-1)y$, proving the formula.

The determinant of complex conjugation on $T$ is $-1$, so its
$+1$ eigenspace has rank one. The map
$(g-1):T^{c=1}\to T$ is injective: a vector in its kernel is fixed
by $c$ and $g$, hence by $h=cgc$ and all of $\Gamma$. Such a vector
is a compatible system of rational $3$-power division points and
is zero because $E(\Q)$ is finite. Thus the displayed quotient
has rank $2-1=1$. No assertion about its finite torsion is needed.
\end{proof}

\begin{proof}[Proof of Theorem~\ref{thm:arbitraryintro}]
Let $N_n=\pi_1((U_n)_{\bar\Q})^{(3)}$ for $n=3^r$.
Lemma~\ref{lem:action} identifies $N_n$ with a $G_\Q$-stable open
subgroup of $N_3$. The stability uses the compatible coordinate
tangents: the local parameter downstairs is the
$3^{r-1}$-st power of the parameter upstairs with leading
coefficient one, so the tangential arithmetic splittings commute
with the covering map. Therefore the action on this actual open
subgroup factors through the same $\Gamma$ as the cubic action.

The image $T'$ of $N_n$ under the compact projection
$N_3\twoheadrightarrow T$ has finite index in $T$. Indeed the
image of an open subgroup under a surjective continuous
homomorphism of profinite groups has finite index. In particular
$T'\otimes\Q_3=T\otimes\Q_3$. By Lemma~\ref{lem:compactrank}
choose $x\in T'$ whose class $\alpha$ in $H^1(\Gamma,T)$ has
infinite order, and choose $a\in N_n$ above $x$.

For each $t\in\Z_3$, use Lemma~\ref{lem:directcocycles} with
$a^{1+3t}$ in place of $a$, and inflate its cocycle along
$G_\Q\twoheadrightarrow\Gamma$. The exponent is well defined
in the procyclic closure of $a$. Its compact class is
\[
\alpha_t=(1+3t)\alpha.
\]
These classes have infinite order and are pairwise distinct.
Inflation $H^1(\Gamma,T)\to H^1(\Q,T)$ is injective: the same
vector exhibiting an inflated class as a coboundary would exhibit
the original class as a coboundary. Therefore the resulting
sections are pairwise geometrically nonconjugate. Their whole
cocycles are unramified outside $3$, since they factor through
$\Gamma$. There are continuum many parameters $t$.

A rational cusp of $X_n$ maps to a rational cusp of $E$. Its
compact class is thus one of the three torsion Kummer classes
in $E(\Q)$, independently of its inertia parameter. Each
$\alpha_t$ is nontorsion, so none of the constructed sections
is globally cuspidal. The compact projection factors through
$N_n^{\ab}$, so this distinction, and pairwise nonconjugacy,
persist in every Galois-stable geometric quotient retaining that
abelianization.

Finally, local point-or-cusp geometry at $3$ would make
$\alpha_t$ a local Kummer class there. Lemma~\ref{lem:cubiclocal}
would then make it a global Kummer class in the finite group
$E(\Q)\simeq\Z/3\Z$, contradicting its infinite order.
Every constructed section therefore fails the Selmer condition
at $3$.
\end{proof}

The choices of an arithmetic basis and of $a$ are noncanonical, but
once made they give the displayed family of full cocycles by a fixed
formula on generators. This construction supplies abstract
noncuspidal sections and proves their local obstruction. It makes
no claim of a noncuspidal Selmer section.

\subsection{A family with zero compact abelianization}

The preceding construction detects noncuspidality through a nontorsion
compact class. A different family shows that the obstruction can be
entirely in the boundary part of geometric homology.

\begin{lemma}[Lifting over a source of $p$-cohomological dimension one]
\label{lem:cdlift}
Let $\Gamma_0$ be profinite with $\operatorname{cd}_p\Gamma_0\leq1$,
and let $N_0$ be a finitely generated pro-$p$ group with a continuous
$\Gamma_0$-action. Every continuous cocycle
$\Gamma_0\to N_0^{\ab}$ lifts to a continuous $N_0$-valued cocycle.
\end{lemma}
\begin{proof}
Pull $N_0\rtimes\Gamma_0\to N_0^{\ab}\rtimes\Gamma_0$ back along
the section specified by the cocycle. The resulting extension
$E\to\Gamma_0$ has kernel $D=[N_0,N_0]$.
Choose a nested cofinal sequence $V_i$ of characteristic open subgroups
of $N_0$. Then $E_i=E/(D\cap V_i)$ has finite $p$-group kernel and
$E=\varprojlim_iE_i$.

Every extension of $\Gamma_0$ by a finite $p$-group $L$ splits.
To see this, induct on $|L|$. The nontrivial characteristic elementary
abelian subgroup $Z=\Omega_1(Z(L))$ is normal in the extension.
First split the quotient extension with kernel $L/Z$ by induction,
then pull back along its section. The remaining extension has kernel
$Z$, and its obstruction lies in $H^2(\Gamma_0,Z)=0$.
The action on $Z$ is the one induced by this extension. Continuous
cohomology applies: an extension with finite kernel has local
continuous sections, by choosing an open subgroup meeting its kernel
trivially, and hence a continuous set-theoretic section on finitely
many clopen translates.

If a section into $E_i$ has been chosen, pull $E_{i+1}\to E_i$ back
along it. The same finite-$p$-kernel argument supplies a compatible
section into $E_{i+1}$. Taking this sequence's inverse limit gives a
continuous section $\Gamma_0\to E$. Its $N_0$-component lifts the
specified abelian cocycle exactly.
\end{proof}

\begin{proposition}\label{prop:boundaryfamily}
For every $n=3^r$, $r\geq1$, there are continuum many pairwise
geometrically nonconjugate full pro-$3$ sections of $U_n/\Q$ whose
whole cocycles are unramified outside $3$ and whose classes in
$H^1(\Q,T_3\Jac(X_n))$ are zero. Each fails point-or-cusp geometry
at $3$. All these assertions persist after passing to any
Galois-stable geometric quotient retaining $N_n^{\ab}$.
\end{proposition}
\begin{proof}
Use $x=X/Z$, $y=Y/Z$ and the standard tangent $t=x$ at
$Q_0=[0:-1:1]$. In its Puiseux fiber functor,
\[
x=t,\qquad y=-(1+t^n)^{1/n}.
\]
The leading constants of $x,y$ are $1,-1$; both have compatible
rational $3$-power roots. The reference section therefore has zero
Kummer classes under these two invertible functions.

Set $M=N_n^{\ab}$. The rational cusps $Q_0$ and
$Q_1=[-1:0:1]$ supply Galois-equivariant maps
$j_0,j_1:\Z_3(1)\to M$ from their inertia groups.
The map induced by the functions $x,y$ satisfies
\[
(q_x,q_y)j_0(z)=(z,0),\qquad
(q_x,q_y)j_1(z)=(0,z),
\]
since these are their orders of zero at the two cusps.
All boundary classes have zero image in $T_3\Jac(X_n)$.

Let $\Gamma=\Gamma_{\{3\}}$ be the arithmetic source above.
The free subgroup $P$ has index two, so
$\operatorname{cd}_3\Gamma=1$: the Hochschild--Serre spectral sequence
and invertibility of two on every $3$-primary module reduce the
assertion to $\operatorname{cd}_3P=1$.
The Kummer cocycle $\kappa_3(3)$ factors through $\Gamma$.
Indeed the fields obtained by adjoining compatible $3$-power roots
of $3$ and roots of unity are unramified outside $3$, and over
$K=\Q(\zeta_3)$ their Galois closures are pro-$3$.
For every $b\in3+9\Z_3$ the abelian cocycle
\[
a_b=(j_0+bj_1)\kappa_3(3)
\]
therefore lives on $\Gamma$. Lemma~\ref{lem:cdlift} lifts it to an
actual $N_n$-cocycle on $\Gamma$. Inflate to $G_\Q$.
The entire cocycle is unramified outside $3$ because this is true
of $G_\Q\to\Gamma$, and its compact abelianization is zero.

Localize at $3$ and apply the valuation to the two coordinate Kummer
classes. The resulting vector is $(1,b)$, so different parameters
give different geometric conjugacy classes.
If a point $(x,y)\in U_n(\Q_3)$ gave this vector, then
$v_3(x)=1$. The equation $y^n=-1-x^n$ would imply $v_3(y)=0$,
contradicting $b\in3+9\Z_3$.
The only rational cusps over $\Q_3$ are $Q_0,Q_1$ and
$Q_\infty=[1:-1:0]$, since $\Q_3$ has no nontrivial
$3$-power roots of unity. Their local coordinate orders give,
respectively, the three lines of possible cusp valuation vectors
\[
(z,0),\qquad(0,z),\qquad(-z,-z)\qquad(z\in\Z_3).
\]
The vector $(1,b)$ is on none of them: $b$ is nonzero and is not one.
Thus the sections are neither point geometric nor cusp geometric
at $3$. This also proves global noncuspidality. The coordinate maps
and all these valuation distinctions factor through $M$, so they
remain valid in every stated quotient.
\end{proof}

The two constructions exhibit different failures of local geometry.
The first is visible in the proper elliptic Tate module; the second
vanishes in compact homology and is detected by two boundary Kummer
coordinates. Neither construction gives a Selmer counterexample.

